\section{Component topology from two weights}
\label{sec:topology-spectra-native}

Report 58's observer is now integrated as
\code{normal\_topology\_spectrum}, after fixing the trace-vocabulary defect
of Section~\ref{sec:incoming-sectors}. It returns the joint distribution of
Euler characteristic, boundary-circle count, orientability and genus or
crosscap number, with binary multiplicities. The earlier coordinate census
can obtain this information by first recovering every distinct component's
full normal vector and then classifying it. The new interface requests only
two additive weights. Its initial integration uses one boundary query and
two surface queries; Section~\ref{sec:weighted-coorientation} subsequently
derives the double histogram when a coorientation certificate succeeds.
This removes the dimension $7t$ coordinate payload when a caller needs
abstract component types.

Polynomial encoded computation of supplied normal-surface topology is
classical~\cite{aht-orbits}. The integration is a smaller observer with
explicit compressed boundary markers, one-sided multiplicity restoration,
and independent source-bound replay. The mathematical construction comes
from report 58; the maintained adaptation adds the explicit monotone-trace
contract and its regression tests. It leaves component coordinates,
boundary slopes and essentiality to their existing observers. Those are
different outputs, so the timing comparison below states its reference's
extra work explicitly.

\subsection{Least representatives from an order-preserving trace}

For an interval relation on $[0,M)\cap\mathbb Z$, a least transversal
contains precisely the minimum point of each orbit. An output interval
$[a,b)$ encodes $b-a$ distinct orbit representatives, not one orbit.
Even $2^{20{,}000}$ isolated points have the one-interval transversal
$[0,2^{20{,}000})$.

The forward compiler maintains an increasing injection from current points
to surviving original labels. Its image is a list of original intervals.
Identity deletion, pairing trimming, periodic merging and transmission
change the generating relation but preserve its equivalence classes on
the same point universe. A suffix truncation discards points equivalent
to retained earlier points, so no unfinished orbit loses its least original
label. A contraction deletes static intervals: each removed point is then
the sole live point of its original orbit, and is emitted. At completion
no live points remain, so the emitted union is exactly the least transversal.
This proof depends on the increasing injection.

The checker independently replays the original orbit proof and reconstructs
the selected set backwards. Undoing a suffix truncation appends unselected
points. Undoing a contraction inserts all its static gaps as selected and
lifts earlier selected points over them. For half-open gaps $[a_i,b_i)$,
write $\ell_i=b_i-a_i$ and $c_i=a_i-\sum_{j<i}\ell_j$. A surviving contracted
point $q$ lifts to
\[
 \lambda(q)=q+\sum_{i:c_i\le q}\ell_i.
\]
Lift a selected interval $[u,v)$ to the hull
$[\lambda(u),\lambda(v-1)+1)$ and include every inserted gap. The hull's
extra points are only those inserted gaps and are already selected, so
the union is exact. Prefix sums and binary searches avoid point expansion.

If the trace has $E$ events and $G$ static-gap intervals, reverse insertion
creates at most $G$ final representative intervals. Forward deletion leaves
at most $1+G$ live original intervals and at most $1+2G$ raw emitted pieces.
Each deleted gap can split at most one live interval; suffix deletion adds
none. During emission, splitting at the $2g$ gap endpoints adds at most
$2g$ pieces, and the live-list sizes telescope over successive contractions.
These observations give the stated bounds.

The list implementation takes
$O(E(1+G)+G\log(G+1))$ endpoint operations for compilation; a conservative
reverse-selector bound is $O((E+G)(G+1)\log(G+2))$, besides ordinary orbit
replay. Endpoint bit lengths are bounded by $\log(M+1)$. Neither output
multiplicity nor a long interval is expanded. These are polynomial local
bounds, rather than an assertion that the list implementation is optimal.

\subsection{The vocabulary guard and the reflection counterexample}

A valid orbit-count proof need not satisfy the monotone-label invariant.
On $[0,3)$ with relation $0\sim1$, the reflected count proof is valid, but
ignoring reflection during reverse selection yields the false set
$\{0,1\}$ in place of $\{0,2\}$. Equal cardinality is insufficient.

The maintained reverse selector now explicitly accepts only contraction,
suffix truncation, deletion, trimming, merging and transmission. Any other
event, including global reflection, rejects the transversal certificate.
The forward compiler likewise reports an unsupported monotone-trace request
when a source-bound count proof contains reflection. Discovery explicitly
requests the forward sweep, so a future default direction change cannot
silently change its proof language.

The guard is on the actual operations, not merely the version number.
For example, a valid version-three count proof containing only supported
monotone events still has the same selector theorem. Two cancelling global
reflections remain outside this compiler's contract, even though the net
map is the identity. This is an explicit language restriction, not a
negative mathematical claim about the source relation. The ordinary forward
AHT algorithm is complete for every supported interval input, so the observer
remains complete on its supplied-vector domain.
Monotone here describes the surviving point labels; reversed arc pairings
and nonorientable surfaces remain fully supported.

Tests replay the retained counterexample with the count checker, reject it
with the corrected selector, and obtain the correct intervals through forward
discovery. Another test places a cancelling reflection pair in a normal
surface's boundary proof: count replay accepts the pair, but the outer topology
checker rejects that unsupported selector language. Producer and checker
continue to propagate caller cancellation, including \code{ValueError}
from a false-valued callable. Independent verification never invokes the
forward selector or discovery.

\subsection{One boundary marker per actual boundary circle}

Build the ordinary interval relation for the surface boundary and compute
its least transversal. Include those boundary-edge points in the full
normal-edge universe, respecting the separate block offsets. Each surface
component then receives exactly one marker for each of its boundary circles.
An additive marker weight therefore measures the actual count $b$, rather
than the number of normal arcs or intersections with a boundary edge.

For the other weight, use the existing normal-cell ownership convention:
one per edge point, minus one per global normal arc including boundary arcs,
plus one per owned
normal disc. Its sum on a component is $\chi$. Thus the full-surface query
has integer weight vector $(\chi,b)$, of dimension two independent of $t$.

In the orientable ambient manifold, doubled normal coordinates form the
orientation cover. Every original point $i$ has adjacent lifts $2i,2i+1$.
Replacing each weight interval $[a,b)$ by $[2a,2b)$ with the same value
lifts both observations. Each boundary circle has two separate lifts:
its tangent line is orientable and its inward collar line is trivial, so
the orientation character restricts trivially there. Consequently the same
original boundary markers correctly count boundary circles of the double;
no second boundary discovery is required.

\subsection{Sparse inversion and complete type reconstruction}

Let $H(v)$ be the base histogram and $K(v)$ the double histogram for
$v=(\chi,b)\in\mathbb Z^2$. Split the base into orientable and nonorientable
counts $O,A$. The geometric cover gives
\[
 H(v)=O(v)+A(v),\qquad K(v)=2O(v)+A(v/2),
\]
where a nonintegral half or an absent base key contributes zero. For $v\ne0$,
process present base keys in increasing $|\chi|+b$ and set
\[
 A(v)=\frac{2H(v)+A(v/2)-K(v)}2,\qquad O(v)=H(v)-A(v).
\]
An integral nonzero half has smaller norm, so it is already known. Only
present keys need processing. At zero the two equations instead give
$A(0)=2H(0)-K(0)$ and $O(0)=K(0)-H(0)$; this distinguishes tori and Klein
bottles with the same zero signature.

Reject nonintegral or negative counts and counts exceeding $H(v)$. Then
reconstruct the entire base and cover histograms, including keys present
only in the cover. This proves existence as well as uniqueness of the
computed split. For instance, three M\"obius bands and five annuli give
$H(0,1)=3$, $H(0,2)=5$ and $K(0,2)=13$. The half-signature contribution
is essential to recover zero nonorientable annuli.

An orientable row has genus $(2-\chi-b)/2$, requiring a nonnegative even
defect. A nonorientable row has $2-\chi-b\ge1$ crosscaps. The independent
algebraic checker reconstructs the complete histograms from proposed rows;
it does not run this inversion. It checks classification, canonical ordering,
strict integral fields and full-support equality. Integral signatures are
required: reducing coordinates modulo two before inversion destroys the
unique-halving argument.

For $s$ sparse keys with fixed dimension, the recurrence uses a linear
number of integer operations after sorting. With a balanced comparison
map and $\ell$-bit fields including temporary sums, the bound is
$O(s\ell\log(s+1))$. The Python implementation uses dictionaries, whose
integer keys can collide adversarially. Its conservative bound is
$O(s^2\ell+s\ell\log(s+1))$. Both are polynomial, but the sharper map
bound is not advertised as the deterministic worst-case cost of these
Python dictionaries.

\subsection{Peeling links and restoring one-sided multiples}

The default observer uses the existing quadrilateral-content reduction.
Peel the minimum triangle coefficient at each global vertex, recording
that many vertex links. Matching triangle differences are quadrilateral
differences. If all quadrilaterals vanish, the peeled core vanishes too.
Otherwise the quadrilateral gcd $g$ divides every remaining triangle,
since one triangle at each vertex is zero and differences propagate over
the connected corner graph. Query the divided core.

For a core orientable component, multiplying normal coordinates by $g$
gives $g$ copies. For a nonorientable component it gives $\lfloor g/2\rfloor$
copies of its connected orientation double and one original precisely when
$g$ is odd. Double both $\chi$ and $b$ on the orientation-double row.
Restore peeled boundary-vertex links as discs and interior-vertex links
as spheres, and combine equal types. A separate checker applies these
typewise equations without calling the restoration producer.

The outer source-bound certificate binds the original labelled triangulation
and coordinate vector, not just the quotient. Replay independently reconstructs
the peel and divisibility, the boundary selector, both weighted histograms,
the cover equations and the original summary. Euler totals must also equal
the original geometry. Reducing core size does not remove input reading,
fingerprinting, gcd, division or serialization costs on the original bits.
The shared cycle allowance covers all three actual orbit discoveries;
compilation and verification retain their cooperative checks. Incomplete
discovery publishes no spectrum or completed certificate.

\subsection{Validation, measurements and composition}

The current-tree audit compares all 1,275 vectors in the maintained
component corpus under both core settings and both merger thresholds:
5,100 complete spectra and independent certificates agree with the frozen
component-level oracle. It retains 56 representative full proofs. A fresh
Regina audit checks 1,395 API results and certificates on 664 bounded vectors
from 26 triangulation isomorphism classes, including genus, crosscap and
boundary multiplicities. The explicit reflection counterexample is rejected
by the maintained checker while its original count proof remains valid.
All 43 focused tests and all 1,312 maintained tests pass. Focused checks cover
source and spectrum mutations, cover-only keys, zero signatures, binary
multiplicities, coordinate reduction, callback propagation, exact cycle caps,
producer-disabled replay and the newly rejected reflection language.

The isolated benchmark has five randomized rounds, with a separate control
for each arm. All 300 measured calls and 60 warm-ups complete. The first
comparison uses the existing coordinate census followed by connected topology
classification as its reference. It computes the same final abstract spectrum,
while additionally producing component normal coordinates. Both arms include
independent proof replay and complete certificate serialization and hashing;
input construction is outside the timer. All spectra agree, and proof records
are deterministic within each arm.
\begin{center}
\small
\begin{tabular}{lrrrrr}
\toprule
Input & coord. ms & spectrum ms & coord./new & coord. A/A & new A/A\\
\midrule
Layered 8 & 8.881 & 9.467 & 0.946 & 0.992 & 0.996\\
Layered 32 & 65.535 & 73.658 & 0.952 & 1.010 & 1.007\\
Layered 64 & 239.303 & 261.880 & 0.916 & 0.958 & 1.006\\
Layered 32, binary scale & 120.869 & 69.048 & 1.727 & 0.970 & 0.995\\
Empty & 0.375 & 0.326 & 1.126 & 1.011 & 0.999\\
Vertex link & 2.406 & 0.381 & 6.412 & 1.001 & 1.028\\
M\"obius band & 1.350 & 1.137 & 1.188 & 1.010 & 1.004\\
Even M\"obius multiple & 1.393 & 1.131 & 1.231 & 0.991 & 1.003\\
Mixed components & 11.517 & 2.250 & 5.153 & 1.004 & 1.008\\
Binary mixture & 15.039 & 2.679 & 5.650 & 1.012 & 0.999\\
Klein bottle & 2.175 & 1.629 & 1.333 & 0.999 & 0.993\\
\bottomrule
\end{tabular}
\end{center}


This comparison is mixed. Vertex links, the ordinary mixture and the binary
mixture improve by factors 6.412, 5.153 and 5.650 respectively. The binary
mixture's proof shrinks from 170,120 to 11,530 bytes. The primitive layered
examples regress by about 5--9 percent despite the smaller weight dimension.
Their connected-coordinate reference benefits from the existing single-orbit
shortcut and the finite coorientation-derived double count. The new observer
still discovers its full weighted double. Dimension reduction alone does
not establish a universal speedup; eliminating that redundant double when
a finite coorientation certificate exists is a further integration candidate.

The binary-scaled layered case multiplies the 32-tetrahedron meridian by
$2^{2048}+1$. The binary mixture uses independently large sphere, vertex-disc
and one-sided populations, including a coefficient exceeding $2^{4096}$.
Its spectrum is restored without expanding those populations.

The second comparison isolates the existing coordinate-core reduction,
using the new direct two-weight observer as the reference and the reduced
two-weight observer as the new arm. Both return the same original spectrum.
\begin{center}
\small
\begin{tabular}{lrrrrr}
\toprule
Input & direct ms & core ms & direct/core & direct A/A & core A/A\\
\midrule
Layered 32, binary scale & 111.527 & 69.691 & 1.618 & 1.007 & 0.998\\
Vertex link & 2.309 & 0.364 & 6.380 & 0.988 & 1.014\\
Even M\"obius multiple & 1.203 & 1.151 & 1.044 & 1.010 & 0.995\\
Binary mixture & 15.049 & 2.674 & 5.616 & 1.000 & 0.999\\
\bottomrule
\end{tabular}
\end{center}


Reduction improves the scaled layered case by 1.618 and the binary mixture
by 5.616. The binary mixture's direct proof has 370,470 bytes, compared with
11,530 after reduction. In the pure-link case it removes all expensive
orbit queries; on an even M\"obius multiple the change is much smaller.
Across both comparisons, control medians range from 0.958 to 1.028. Factors
are medians of paired ratios, not ratios of the displayed timing medians.
These are supplied-normal-vector API measurements, rather than complete
knot-recognition timings or a complexity fit.


The runtime lives in the six new \code{fastunknot} observer modules and has
an explicit \code{census}/\code{verify} command-line interface through
\code{python -m fastunknot.normal\_topology}. The preserved report remains
unchanged. The maintained driver
\path{../fast/normal_orbit_research/spectra.py} has \code{audit} and
\code{benchmark} modes with an output path and optional rounds. Fresh
records, proofs, source pins, tests and publication checks are
\path{data/topology-spectra-*}.

If an outer algorithm makes $Q(n)$ queries of encoded size at most $S(n)$,
this observer contributes $Q(n)\operatorname{poly}(S(n))$ time. Controlling
both functions at quasipolynomial size would preserve that class for this
stage. The integration does not prove those hypotheses. Abstract component
types alone do not supply a useful vector, essential boundary slope,
component coordinates, cutting instructions or attachment maps. The sector
search and existing diagram-bound positive certificates remain relevant
to those separate obligations.
